\begin{abstract}
Vision-language models describe social scenes well but often choose the wrong action in them. We study how to improve norm-grounded action selection, and what that selection actually rests on, through five interventions and analyses evaluated mainly on EgoNormia, each against matched controls. Targeted internal visual regularization of Qwen3.5-9B lowers vote negative log-likelihood on NormLens beyond supervised fine-tuning and a matched random perturbation, although only one training seed has been evaluated. Reinforcement learning on descriptive movie-clip reasoning improves in-domain accuracy, but its transfer to EgoNormia is uncertain ($+2.5$ points joint accuracy), smaller than answer-only fine-tuning ($+4.3$) and smaller than the effect of a longer response budget ($+7$). Projecting an appearance-trust direction out of Qwen3-VL's activations cuts cue-induced answer changes from 24.5\% to 17.8\% without reducing responsiveness to behavior, about as well as a constitution prompt; synthetic faces carry no measurable trust effect. In Qwen3-Omni, audio's apparent gain over video comes mostly from fewer abstentions, and vanishes once the ``None'' option is removed. Explicit depth, segmentation and pose maps, as well as spatial fine-tuning, give no gain over matched controls (pooled $+0.2$ points), and spatial fine-tuning lowers accuracy. Across ideas, many apparent gains reflect abstention and output format rather than better normative reasoning, which makes such controls essential.
\end{abstract}
